\section{\method{}: Calculating the exact influence function at scale} We now
present \method{}, our method for nearly-optimal single-model data attribution.
The skeleton of our method is conceptually straightforward: we exactly calculate
the influence function in large-scale learning settings.
In this section, we first formally define the influence function;
we then define a specific class of learning algorithms that we will consider;
and finally, we present our method for calculating the exact influence function
for this class of learning algorithms.

\subsection{The influence function} %
\label{sec:influence_fn}
At the core of our method is a statistical primitive
known as the {\em influence function approximation}
\citep{hampel1947influence,koh2017understanding,giordano2019swiss}.
The main idea is to approximate the model output $f$ for a given
data weighting
$\wvec$ using the following first-order Taylor expansion:
\begin{equation}
  \label{eq:infl_fn}
  \hat{f}(\wvec) := f(\mathbf{1}_n) +
  \left(\frac{\partial f(\wvec)}{\partial
  \wvec}\bigg|_{\wvec=\mathbf{1}_n}\right)^\top  (\wvec -
  \mathbf{1}_n);
\end{equation}
The key term in this estimate is the gradient
$\nicefrac{\partial f(\wvec)}{\partial \wvec}$ evaluated at
$\wvec=\mathbf{1}_n$, called the \textit{influence function}.
Intuitively, this term captures the effect of infinitesimally up- or
down-weighting each training example on the model output.
While well-defined, this quantity is not trivial to compute:
after all, it is the gradient ``through'' the process of training
the model with algorithm $\alg{}$.



\paragraph{The main challenge of influence functions in large-scale settings:
computing the gradient.} When the learning algorithm $\alg{}$ is a {\em convex}
optimization algorithm (e.g., linear regression, logistic regression, etc.),
the influence function is straightforward to compute. Indeed, in such settings,
the gradient $\nicefrac{\partial f(\wvec)}{\partial \wvec}$ has a simple closed
form (via implicit differentiation),
and the first-order Taylor expansion \eqref{eq:infl_fn} yields
near-perfect estimates of the model output $f(\wvec)$ will behave on many
choices of new data weightings $\wvec$
\citep{koh2017understanding,rad2018scalable,koh2019accuracy,giordano2019swiss}.

In large-scale, non-convex settings (e.g., in deep learning),
however, the influence function is difficult to compute:
no such closed form exists.
Instead, large-scale data attribution methods that are based on
\eqref{eq:infl_fn}
must {\em approximate} the influence function
\citep{koh2017understanding,bae2022if,park2023trak,bae2024training}.
And while these methods have shown promise,
they are not nearly as effective as the influence function is in
analogous convex
settings.





\subsection{Focus: iterative smooth learning algorithms}
Before describing our procedure, we first formalize the class of
learning algorithms
$\alg{}$ that we will consider, namely ones that are {\em iterative}
and {\em smooth}.
By restricting our focus to this class
of learning algorithms, we ensure that the influence function is well-defined.
Note: this class of algorithms is extremely general and captures,
e.g., large-scale (transformer-based) language model training or deep image
classifier training (sometimes with slight changes from standard
learning algorithms).

\subsubsection{Iterative learning algorithms}
First, we require that the learning algorithm $\alg{}$
is {\em iterative}, i.e., it takes the form
\begin{align}
  \label{eq:iterative_algo}
  \alg(\wvec) := \st{T} \quad \text{for} \quad
  \st{t+1} &:= h_t(\st{t}, \mathbf{g}_t(\st{t}, \wvec))
  \quad \text{ and } \quad \mathbf{g}_t(\st{t}, \wvec) :=
  \sum_{i\in B_t} w_i \cdot \nabla_{\st{t}} \ell(z_i; \st{t}).
\end{align}
Above,
$\st{t}$ is the optimizer state (including model parameters),
which is iteratively updated by a function $h_t$
starting from a fixed initial state $\st{0}$.
We let the number of training steps be $T$;
$B_t \subset [N]$ is a minibatch sampled at step $t$;
and $\ell(z_i; \st{t})$ is the loss on
sample $z_i$ given optimizer state $\st{t}$.

The vast majority of large-scale learning algorithms are iterative
in this sense---we give a few examples below.
Recall that the learning algorithm $\alg{}$ takes as input a
data weighting $\wvec$ over the training set, and outputs
a machine learning model trained on the weighted dataset.
The algorithm thus encapsulates all aspects of the training setup
beyond the training data weights, including the model architecture,
optimizer, and hyperparameters.

\begin{example}[Training an ResNet with SGD]
  Here, the optimizer state $\st{t}$ is the parameter vector
  $\theta_t$ at step $t$,
  the loss function $\ell(z_i; \st{t})$ is the cross-entropy loss of a ResNet
  with parameters $\st{t}$ on a given training example $z_i$,
  and the update function $h_t$ is the SGD update step
  $$h(\st{t}, \mathbf{g}_t) := \st{t} - \eta_t \cdot \mathbf{g}_t,$$
  where $\eta_t$ is the learning rate at step $t$.
\end{example}

\noindent More complex optimizers can be handled by extending the definition
of the optimizer state $\st{t}$ beyond just the parameter vector,
as we show in the following example.
\begin{example}[Training a language model with Adam]
  Here, the optimizer state $\st{t} = (\theta_t, m_t, v_t)$,
  where $\theta_t$ is the parameter vector at step $t$,
  $m_t$ is the first moment estimate of the gradient,
  and $v_t$ is the second moment estimate of the gradient.
  The loss function $\ell(z_i; \st{t})$ is the cross-entropy loss of a language
  model with parameters $\theta_t$ on a given training example $z_i$,
  and the update function $h_t$ is the Adam update step:
  $$h(\st{t}, \mathbf{g}_t) := \left[
    \begin{array}{c}
      \theta_t - \eta_t \cdot \frac{\sqrt{v_t}}{\sqrt{m_t +
        \epsilon_{\mathrm{root}}} +
      \epsilon} \cdot \mathbf{g}_t \\
      \beta_1 \cdot m_t + (1 - \beta_1) \cdot \mathbf{g}_t \\
      \beta_2 \cdot v_t + (1 - \beta_2) \cdot \mathbf{g}_t^2
    \end{array}
  \right]$$
  where $\eta_t$ is the learning rate at step $t$.
\end{example}



\subsubsection{Smooth learning algorithms}
Finally, for predictive data attribution to even be possible (recalling
Definition \ref{def:predictive_data_attribution}), we must ensure that the
learning algorithm $\alg{}$ is (qualitatively) well-behaved as a function of the
data weights $\wvec$.
Indeed, when this is not the case,
it is unlikely that {\em any} simple predictor $\hat{f}$
will be able to accurately predict $f(\wvec)$ from $\wvec$.
To formalize this requirement, we consider learning algorithms $\alg{}$ that are
{\em smooth} in $\wvec$, as described by \citet{engstrom2025optimizing}.
In particular, a learning algorithm $\alg{}$ is smooth in $\wvec$ if,
for any measurement function $\phi$, small perturbations to the data weights
$\wvec$ result in only small changes to the gradient
$\nicefrac{\partial f(\wvec)}{\partial \wvec}$.

To see why smoothness is necessary to predict $f$ with a simple function, we walk
through a simple example (visualized in Figure \ref{fig:smoothness}).
For a model output function $f$, consider slightly upweighting the
$i$-th training sample,
and measuring the change in the model output $\Delta(\varepsilon) :=
f(\wvec + \varepsilon \mathbf{1}_i) - f(\wvec)$.
If the learning algorithm is smooth, then this change is well-behaved as a
function of $\varepsilon$.
In particular, the change $\Delta(2\varepsilon)$ should be
reasonably approximated by $2\Delta(\varepsilon)$.
On the other hand, if the learning algorithm is \textit{not} smooth,
$\Delta(\varepsilon)$ may change wildly as $\varepsilon$ varies,
precluding any simple prediction method from being able to accurately
predict $f(\wvec)$ from $\wvec$.


\begin{figure}[h]
  \centering
  \begin{tikzpicture}
  \begin{groupplot}[
      group style={
        group size=2 by 1, %
        horizontal sep=1.5cm %
      },
      width=8cm,       %
      height=4cm,      %
      xmin=0, xmax=0.5,
      ymin=-10, ymax=15,
      axis lines=middle, %
      xtick=\empty,      %
      ytick=\empty,      %
      grid=none,         %
      enlargelimits=false, %
      xlabel={$\varepsilon$}, %
      ylabel={{\small $\Delta(\varepsilon) \coloneqq f(\wvec +
          \varepsilon \mathbf{1}_i) -
      f(\wvec)$}} %
    ]

    \nextgroupplot %
    \addplot[
      domain=0:0.46,
      samples=300, %
      blue %
    ]
    { 9 * sin(deg(90*x)) * sin(deg(20*x)) };

    \nextgroupplot %
    \addplot[
      domain=0:.46,
      samples=100, %
      smooth,      %
      red          %
    ]
    { 40*x - 40*x^2 };

  \end{groupplot}
\end{tikzpicture}

  \caption{Smoothness aids predictive data attribution. We plot the change in
    data weights $\varepsilon$ against the change in model output
    $\Delta(\varepsilon)$ for two hypothetical learning algorithms. On the left
    is a non-smooth setting where the
    gradient $\nicefrac{f(\wvec)}{\wvec}$ varies wildly with
    $\varepsilon$. On the
  right is a smooth setting where the change is well-behaved. }
  \label{fig:smoothness}
\end{figure}

\begin{remark}[How restrictive is smoothness?]
  Unlike iterativity, smoothness is not an inherent
  property of standard learning algorithms.
  Indeed, many standard training routines are \underline{not}
  smooth.
  In such cases, there is
  \underline{\smash{no}} good data attribution method $\hat{f}$
  satisfying a natural ``additivity'' property \citep{saunshi2023understanding}
  (ruling~out essentially all known data attribution methods).
  Fortunately, however, as observed by prior work,
  one can often construct a ``smooth counterpart''
  to any given non-smooth learning algorithm \citep{engstrom2025optimizing}.
  Motivated by this finding (and the impossibility above),
  we focus our attention on smooth learning algorithms.
\end{remark}


\subsection{Calculating the exact influence function}
To compute the exact influence function for a model output function $f$
that is the output of an iterative smooth learning algorithm $\alg{}$,
we leverage recent developments in {\em metagradient} calculation
\citep{maclaurin2015gradient,franceschi2017forward,lorraine2020optimizing,engstrom2025optimizing}.
A metagradient is a gradient of a machine learning model's output
with respect to a design choice made prior to training.
In recent work, \citet{engstrom2025optimizing} present an algorithm
called \replay{}
that {\em exactly} calculates the metagradient for iterative and smooth
learning algorithms.

\paragraph{Calculating the influence function with \replay{}.}
Observe that when the design choice is the data weighting $\wvec$,
the metagradient---the gradient of the model output $f$ with respect to
$\wvec$---is precisely the influence function.
We can thus directly apply the \replay{} algorithm
\citep{engstrom2025optimizing}
to calculate the exact
influence function.
Adapted to our setting, \replay{} calculates the metagradient by
exploiting the following identity, which follows from the chain rule
applied to the computation graph illustrated in Figure
\ref{fig:computation_graph}:
\begin{align}
  \nabla_{\wvec} f(\wvec)
  &= \sum_{t=0}^{T-1} \underbrace{
    \frac{\partial f(\wvec)}{\partial \st{t+1}} \cdot
    \frac{\partial h_t(\st{t}, \mathbf{g}_t(\st{t}, \wvec))}{\partial \wvec}
  }_{\text{contribution of $\wvec$ to $f(\wvec)$ through $\st{t+1}$}}
  \label{eq:replay_meta_grad}
\end{align}

\begin{figure}[t]
  \centering
  \input{figures/computation_graph}
  \caption{Forward computation graph for a model output function $f$
    mapping from data weights $\wvec$ to the model output.
    The exact influence function $\nicefrac{\partial f(\wvec)}{\partial \wvec}$
    is the {\em metagradient} of the model output with respect to the
  data weights $\wvec$.}
  \label{fig:computation_graph}
\end{figure}

Motivated by this observation, \replay{} operates as follows:
\begin{enumerate}
  \item Initialize $\Delta_T = \frac{\partial f(\wvec)}{\partial
    \st{T}} = \nabla_{\st{T}} \phi(\st{T})$
  \item For each $t = T-1, \ldots, 0$:
    \begin{enumerate}
      \item Load state $\st{t}$ and minibatch $B_t$
      \item Calculate $\beta_t = \nabla_{\wvec} \left(h_t(\st{t},
        \mathbf{g}_t(\st{t}, \wvec))^\top \Delta_{t+1}\right)$,
        the contribution of $\wvec$ to $\phi(\st{T})$ through the
        $t$-th step of the learning algorithm
      \item Advance $\Delta_t = \nabla_{\st{t}} \left(h_t(\st{t},
        \mathbf{g}_t(\st{t}, \wvec))^\top \Delta_{t+1}\right)$,
        which is $\nicefrac{\partial f(\wvec)}{\partial \st{t}}$
    \end{enumerate}
  \item Return $\beta = \sum_{t=0}^{T-1} \beta_t$ as the exact
    influence function
\end{enumerate}

\noindent By leveraging an efficient data structure to
load the states and minibatches,
\replay{} is able to calculate the exact influence function
for IDS model output functions at a computational cost of
$T + T\log(T)$ total training steps, and $\log(T)$ memory.
We refer the reader to \citet{engstrom2025optimizing} for a complete
description of the algorithm.








